% =====================================================================
%  POSTER A0 (vertical)  -  PFC I  -  UTEC
%  Tema NARANJA · sin logos · tipografía Avenir Next
%  Compilar con XeLaTeX / Tectonic:   tectonic poster.tex
% =====================================================================
\documentclass[final]{beamer}
\usepackage[size=a0, scale=1.0, orientation=portrait]{beamerposter}

% --- Idioma ---------------------------------------------------------
\usepackage{polyglossia}
\setdefaultlanguage{spanish}
\setotherlanguage{english}

% --- Tipografía (Avenir Next → Helvetica Neue → fallback) -----------
\usepackage{fontspec}
\IfFontExistsTF{Avenir Next}{%
  \setsansfont{Avenir Next}[UprightFont=* Regular, BoldFont=* Bold,
    ItalicFont=* Italic, BoldItalicFont=* Bold Italic]%
}{%
  \IfFontExistsTF{Helvetica Neue}{\setsansfont{Helvetica Neue}}{}%
}
\IfFontExistsTF{Avenir Next Heavy}{%
  \newcommand{\titlefont}{\fontspec{Avenir Next Heavy}}%
}{%
  \IfFontExistsTF{Avenir Next Black}{\newcommand{\titlefont}{\fontspec{Avenir Next Black}}}{\newcommand{\titlefont}{\bfseries}}%
}
% Fuente para títulos de sección (Futura: distinta del cuerpo y más ligera)
\IfFontExistsTF{Futura}{\newcommand{\sectfont}{\fontspec{Futura}}}{\newcommand{\sectfont}{\sffamily}}
\renewcommand{\familydefault}{\sfdefault}

% --- Gráficos y layout ---------------------------------------------
\usepackage{graphicx}
\usepackage{tikz}
\usetikzlibrary{positioning, arrows.meta, calc, shapes.geometric, fit, backgrounds}
\usepackage[most]{tcolorbox}
\usepackage{booktabs}
\usepackage{array}
\usepackage{colortbl}
\usepackage{ragged2e}
\usepackage{amsmath, amssymb}

\usetheme{default}
\setbeamertemplate{navigation symbols}{}
\setbeamertemplate{itemize item}{\color{uorange}\textbf{\textbullet}}
\setbeamertemplate{itemize subitem}{\color{uorange2}\(\circ\)}
\setbeamercolor{itemize item}{fg=uorange}
\setbeamercolor{normal text}{fg=ink}
\setbeamercolor{structure}{fg=uorange}

% --- Paleta UTEC (teal institucional, extraída del logo #107688) ----
\definecolor{uorange}{HTML}{107688}   % primario: teal UTEC
\definecolor{uorange2}{HTML}{1A93A6}  % teal medio/brillante (realces)
\definecolor{udark}{HTML}{074A55}     % teal profundo (header / títulos)
\definecolor{utan}{HTML}{E8F4F6}      % tinte claro teal (paneles)
\definecolor{uborder}{HTML}{BFDCE1}   % borde teal sutil
\definecolor{ink}{HTML}{1B2A30}       % texto
\definecolor{muted}{HTML}{5C6B73}     % texto secundario
\definecolor{upos}{HTML}{2E7D5B}      % positivo
\definecolor{uneg}{HTML}{B23A2E}      % negativo / MTD

\graphicspath{{images/}}

% --- Badge numerado para los títulos de sección ---------------------
\newcommand{\badge}[1]{\tikz[baseline=-0.7ex]\node[circle, fill=white,
  text=uorange, font=\normalsize\bfseries\sectfont, inner sep=0pt,
  minimum size=11mm]{#1};}

% --- Caja de bloque estándar (con badge) ----------------------------
\newtcolorbox{pblock}[2]{%
  enhanced, breakable,
  colback=white, colframe=uborder,
  colbacktitle=uorange, coltitle=white,
  fonttitle=\Large\sectfont,
  title={\badge{#1}\hspace{3mm}#2},
  titlerule=0.6mm, titlerule style={udark},
  boxrule=0.3mm, arc=3mm,
  left=8mm, right=8mm, top=6mm, bottom=6mm,
  toptitle=3mm, bottomtitle=3mm,
  fontupper=\normalsize\sffamily,
  drop fuzzy shadow=black!12,
}

% --- Caja "destacado" para resultados clave -------------------------
\newtcolorbox{hlbox}{%
  enhanced,
  colback=uorange!12, colframe=uorange,
  boxrule=0.4mm, arc=2mm,
  left=6mm, right=6mm, top=3mm, bottom=3mm,
  fontupper=\large\sffamily\bfseries,
}

% --- Caja secundaria (pregunta, etc.) -------------------------------
\newtcolorbox{softbox}{%
  enhanced, colback=uorange!7, colframe=uorange,
  boxrule=0.3mm, arc=2mm,
  left=5mm, right=5mm, top=3mm, bottom=3mm,
  fontupper=\normalsize\sffamily,
}

% --- Items compactos ------------------------------------------------
\newcommand{\bitem}{\textcolor{uorange}{\textbullet}\,}

% --- Figura enmarcada (marco teal fuerte para que destaque) ----------
\newcommand{\pfig}[3][0.99\linewidth]{%
  \begin{tcolorbox}[enhanced, colback=utan, colframe=uorange, boxrule=0.4mm,
     arc=2.5mm, left=4mm, right=4mm, top=4mm, bottom=3mm,
     drop fuzzy shadow=black!14, borderline north={2.2mm}{0mm}{uorange2}]
    \centering\includegraphics[width=#1]{#2}\\[2mm]
    {\small\color{udark}\sffamily #3}
  \end{tcolorbox}}

% =====================================================================
\begin{document}
\begin{frame}[t]
% --------------------------- ENCABEZADO ------------------------------
% Banda naranja con degradado, SIN logo, texto blanco.
\begin{tcolorbox}[enhanced, boxrule=0mm, arc=0mm,
  colback=uorange, colframe=uorange,
  borderline south={5mm}{0mm}{udark},
  left=12mm, right=12mm, top=9mm, bottom=9mm,
  drop fuzzy shadow=black!30]
\centering
{\bfseries\color{white}\huge Mitigación de la resistencia evolutiva en glioblastoma\\[1mm]
 mediante aprendizaje por refuerzo multiagente adversarial}\\[6mm]
{\large\color{white!92} Marco \textbf{GBMARL}: terapia y tumor co-evolucionan por \emph{self-play} MAPPO-CTDE sobre un simulador calibrado con datos reales}\\[7mm]
{\large\color{white!95} Kalos B.\ Lazo Mera \quad\textbullet\quad Gianpier A.\ Segovia Ureta \quad\textbullet\quad Asesor: V.\ E.\ Martínez Abaunza}\\[2mm]
{\color{white!85} Universidad de Ingeniería y Tecnología (UTEC) \,\(\cdot\)\, Ciencia de la Computación \,\(\cdot\)\, PFC I \,\(\cdot\)\, 2026}
\end{tcolorbox}

\vspace{5mm}

% ============================ CUERPO =================================
\begin{columns}[t]

% --------------------- COLUMNA IZQUIERDA ----------------------------
\begin{column}{0.485\textwidth}

% --- 1. PROBLEMA ---
\begin{pblock}{1}{Problema y contexto clínico}
\textbf{\color{udark}Glioblastoma multiforme (GBM):} tumor cerebral primario maligno más frecuente y agresivo en adultos.
\begin{itemize}\itemsep2pt
  \item Mediana de supervivencia \textbf{12--15 meses}; supervivencia a 5 años \textbf{$<$7\,\%}.
  \item Estándar (protocolo de Stupp): cirugía + radioterapia + \textbf{temozolomida} (TMZ). Casi todos los pacientes recaen con \emph{resistencia a la TMZ}.
\end{itemize}

\textbf{\color{udark}La recaída no es accidental, es evolutiva.} La dosis máxima tolerada (MTD) ejerce una presión selectiva que \textbf{expande los subclones resistentes}: al eliminar las células sensibles, libera a las resistentes de la competencia.

\vspace{2mm}
\textbf{\color{uneg}$\blacktriangle$\ Brecha que aborda este trabajo}
\begin{itemize}\itemsep2pt
  \item \textbf{RL de agente único:} trata al tumor como una dinámica \emph{fija}; no anticipa ni responde estratégicamente.
  \item \textbf{Heurísticas adaptativas (Gatenby):} reglas fijas encendido/apagado, \emph{no optimizadas}.
  \item \textbf{Modelos teórico-de-juegos:} resuelven al tumor analíticamente $\Rightarrow$ pobreza de comportamientos adversariales.
  \item \textbf{Métricas de carga aislada:} premian la reducción agresiva del tumor $\Rightarrow$ \emph{inducen} resistencia (engañosas).
\end{itemize}

\vspace{2mm}
\begin{softbox}
\textbf{\color{udark}Pregunta de investigación.} ¿Modelar el GBM como un \emph{juego adversarial de dos agentes} (terapia \emph{vs.}\ tumor), resuelto por MARL sobre un simulador calibrado con datos reales, permite descubrir políticas de dosificación que \textbf{retrasen la dominancia resistente} frente a MTD y a la heurística de Gatenby? ¿Aporta el crítico centralizado (CTDE) sobre el aprendizaje independiente (IPPO)?
\end{softbox}
\end{pblock}

\vspace{23mm}

% --- 2. OBJETIVOS ---
\begin{pblock}{2}{Objetivos}
\textbf{\color{udark}General.} Diseñar \textbf{GBMARL}, un marco de aprendizaje por refuerzo multiagente adversarial que modele la terapia y la resistencia del glioblastoma como agentes en competencia sobre un simulador calibrado con datos reales, y evaluar su capacidad de \textbf{retrasar la intratabilidad} frente a líneas base clínicas.

\vspace{2mm}
\textbf{\color{udark}Específicos.}
\begin{enumerate}\itemsep1pt
  \item Construir y calibrar un \textbf{simulador de EDO} (competencia Lotka--Volterra + farmacodinámica tipo Hill) con datos reales (GDSC2 \& DepMap).
  \item Formular el \textbf{entorno multiagente} con observabilidad parcial clínica e implementar \textbf{MAPPO-CTDE} y su ablación \textbf{IPPO}.
  \item Entrenar y validar con \textbf{múltiples semillas}, reportando significancia estadística e \textbf{interpretando el mecanismo} de la política.
  \item Aislar la contribución de CTDE (\textbf{ablación}) y evaluar la \textbf{sensibilidad} a las constantes de calibración.
\end{enumerate}
\end{pblock}

\vspace{23mm}

% --- 3. METODOLOGIA ---
\begin{pblock}{3}{Metodología --- el marco GBMARL}
\textbf{\color{udark}Juego adversarial de suma no nula.} Dos agentes en \emph{self-play} sobre un simulador de EDO:
\begin{itemize}\itemsep1pt
  \item \textbf{Agente de terapia} elige la dosis $u(t)\in[0,u_{\max}]$ $\to$ maximizar los días con tumor manejable.
  \item \textbf{Agente tumoral} elige la transición fenotípica $\phi\in[0,\phi_{\max}]$ $\to$ forzar la dominancia resistente.
\end{itemize}

\textbf{\color{udark}Simulador calibrado} (Runge--Kutta 4):
\[
\small
\begin{aligned}
\tfrac{dS}{dt} &= \alpha_S S\!\left(1-\tfrac{S+R}{K}\right)-\delta_S(c)S-\phi S\\
\tfrac{dR}{dt} &= \alpha_R R\!\left(1-\tfrac{S+R}{K}\right)-\delta_R(c)R+\phi S\\
\tfrac{dc}{dt} &= -\lambda_c c+u(t),\qquad \delta_X(c)=\delta_X^{\max}\tfrac{c^{h}}{\mathrm{IC50}_X^{h}+c^{h}}
\end{aligned}
\]
\bitem \textbf{Calibración:} datos GDSC2+DepMap fijan la \emph{potencia} del fármaco; la literatura ancla el techo de muerte. $\mathrm{IC50}_S{=}0.36$, $\mathrm{IC50}_R{=}3.27$ $\Rightarrow$ \textbf{brecha $\approx 9\times$} (34 líneas recuperadas).

\bitem \textbf{Observabilidad parcial clínica:} la terapia observa $(S{+}R,\,c)$ pero \textbf{no} la fracción resistente $\Rightarrow$ justifica un \textbf{crítico centralizado} $V(s^{\text{global}})$.

\bitem \textbf{MAPPO-CTDE:} actor local + crítico con estado conjunto; ventajas GAE; \emph{self-play} adversarial.

\bitem \textbf{Métrica honesta (TTP-combinado):} días hasta fallar \emph{control} de carga \textbf{o} \emph{tratable} (mayoría resistente), reportando el \textbf{modo de falla}.

\vspace{2mm}
\begin{columns}[T,onlytextwidth]
\begin{column}{0.49\textwidth}
\pfig{images/pipeline_diagram.png}{Pipeline reproducible (\emph{inside-out}): cada fase de-riskea la siguiente.}
\end{column}
\begin{column}{0.49\textwidth}
\pfig{images/dynamics_baseline.png}{Dinámica del simulador calibrado: poblaciones S/R según régimen de dosis.}
\end{column}
\end{columns}
\end{pblock}

\end{column}

% --------------------- COLUMNA DERECHA ------------------------------
\begin{column}{0.485\textwidth}

% --- 4. RESULTADOS ---
\begin{pblock}{4}{Resultados preliminares}
\textbf{\color{udark}Comparación bajo la métrica TTP-combinado} (tiempo hasta la intratabilidad):
\vspace{1mm}
\begin{center}
\renewcommand{\arraystretch}{1.3}
\begin{tabular}{@{}lcc@{}}
\toprule
\rowcolor{uorange}{\color{white}\textbf{Estrategia}} & {\color{white}\textbf{TTP (días)}} & {\color{white}\textbf{Modo de falla}}\\
\midrule
Sin tratamiento        & 12 & carga progresó\\
MTD (dosis máx.)       & 13 & resistencia mayoría\\
Gatenby (adaptativa)   & 27 & resistencia mayoría\\
\rowcolor{uorange!22}\textbf{MAPPO (aprendida)} & \textbf{33{,}5 $\pm$ 5{,}8}\;\,(mejor 41) & resistencia mayoría\\
\bottomrule
\end{tabular}
\end{center}
{\small Validación multi-semilla ($n{=}15$): \textbf{$p<0{,}001$} vs.\ MTD y vs.\ Gatenby. Corrida completa reproducible (120\,000 pasos, $n{=}5$): MAPPO supera a Gatenby en \textbf{5/5} semillas.}

\vspace{2mm}
\begin{hlbox}
\centering La política aprendida \textbf{retrasa la intratabilidad} $\sim$2{,}6$\times$ frente a MTD y $+$24\,\% (media) / $+$52\,\% (mejor corrida) sobre Gatenby.
\end{hlbox}

\vspace{2mm}
\pfig{images/evaluation_ttp.png}{\textbf{Izq.:} carga tumoral. \textbf{Der.:} fracción resistente. La MTD controla la carga pero lleva la resistencia al 100\,\%; la política aprendida la retrasa.}
\end{pblock}

\vspace{13mm}

% --- 5. MECANISMO + ABLACION ---
\begin{pblock}{5}{Mecanismo descubierto y ablación CTDE}
\begin{columns}[T,onlytextwidth]
\begin{column}{0.5\textwidth}
\textbf{\color{udark}Dosificación \emph{pulsada}.}\\[1mm]
El agente \textbf{redescubre sin ser programado} el principio de la terapia adaptativa: dosis \textbf{intermitente} que preserva una población \textbf{sensible competidora}, frenando la expansión de la resistente.
\pfig[0.99\linewidth]{images/policy_seed0.png}{Poblaciones, fracción resistente y acción de dosis: dosificación pulsada.}
\end{column}
\begin{column}{0.5\textwidth}
\textbf{\color{udark}Ablación CTDE vs.\ IPPO.}\\[1mm]
\small
\begin{tabular}{@{}lc@{}}
MAPPO (CTDE) & $34{,}8\pm5{,}7$ d\\
IPPO (local) & $30{,}2\pm2{,}9$ d\\
\end{tabular}\\[1mm]
MAPPO falla por \emph{resistencia mayoría}; IPPO por \emph{carga} $\Rightarrow$ CTDE sostiene el control más tiempo. Ventaja \textbf{modesta}, caracterizada con honestidad.
\pfig[0.92\linewidth]{images/ablation_ctde.png}{Ablación CTDE: MAPPO (crítico centralizado) vs.\ IPPO (local).}
\end{column}
\end{columns}
\vspace{1mm}
{\small\color{udark}\textbf{Sensibilidad:} al variar $\pm$30\,\% las 6 constantes clave, el orden \textbf{MAPPO $>$ Gatenby $>$ MTD} \emph{nunca} se invierte.}
\end{pblock}

\vspace{13mm}

% --- 6. CONTRIBUCION + ALCANCE ---
\begin{pblock}{6}{Contribución, alcance y limitaciones}
\textbf{\color{udark}Contribuciones:} (1) primer marco MARL adversarial \textbf{calibrado} para GBM; (2) política interpretable (pulsada) que retrasa la resistencia de forma significativa; (3) métrica de evaluación clínicamente honesta; (4) ablación CTDE vs.\ IPPO en oncología.

\vspace{1mm}
\textbf{\color{uneg}Limitaciones (declaradas con honestidad):}
\begin{itemize}\itemsep1pt
  \item Validación \emph{in silico} sobre simulador calibrado \emph{in vitro} (no clínica). GBM + TMZ, modelo \textbf{no espacial} de 2 poblaciones.
  \item Objetivo: \textbf{retrasar} la intratabilidad, \emph{no curar} (la resistencia es inevitable bajo calibración realista).
  \item El \emph{self-play} produce resultados \textbf{bimodales} $\Rightarrow$ se reporta por \textbf{tasa de éxito}, no por promedios engañosos.
\end{itemize}
\end{pblock}

\vspace{13mm}

% --- 7. PLAN DE TRABAJO PFC II (timeline) ---
\begin{pblock}{7}{Plan de trabajo para PFC II}
\centering
\begin{tikzpicture}[
  every node/.style={font=\sffamily},
  dot/.style={circle, fill=uorange, text=white, font=\Large\bfseries\sffamily,
              minimum size=15mm, inner sep=0pt},
  card/.style={anchor=west, text width=0.78\linewidth, align=left, inner sep=1.5mm}]
\draw[uorange, line width=1.4mm, line cap=round] (0,0.7) -- (0,-9.7);
\node[dot] at (0,0) {1};
\node[card] at (1.4,0) {\textbf{\color{udark}Adversario evolutivo más fuerte}\\[0.5mm]\small Ampliar el espacio de acción del tumor ($\phi_{\max}$): amenaza realista y menos bimodalidad.};
\node[dot] at (0,-3) {2};
\node[card] at (1.4,-3) {\textbf{\color{udark}Calibración con datos clínicos}\\[0.5mm]\small Anclar parámetros a imágenes y series temporales de pacientes, no solo a ensayos \emph{in vitro}.};
\node[dot] at (0,-6) {3};
\node[card] at (1.4,-6) {\textbf{\color{udark}Tratamientos combinados y espacialidad}\\[0.5mm]\small Multi-fármaco (TMZ + agentes) y modelos espaciales de tumoresfera.};
\node[dot] at (0,-9) {4};
\node[card] at (1.4,-9) {\textbf{\color{udark}Traslado a otras histologías y cómputo}\\[0.5mm]\small Extender el marco a otros tumores con resistencia evolutiva; acelerar en GPU/vectorización.};
\end{tikzpicture}
\end{pblock}

\end{column}
\end{columns}

% --------------------------- PIE -------------------------------------
\vspace{4mm}
\begin{tcolorbox}[enhanced, colback=utan, colframe=uborder, boxrule=0.3mm, arc=1.5mm,
  left=6mm, right=6mm, top=2mm, bottom=2mm]
{\small\color{ink}
\textbf{\color{udark}Referencias clave:} Stupp (2005); Gatenby \& Gallaher (2017); Stankova (2019); Pinto et al.\ RARL (2017); Gallagher (2024); Yu et al.\ MAPPO (2022); de Witt et al.\ IPPO (2020); GDSC2; DepMap.}
\end{tcolorbox}

\end{frame}
\end{document}
